\clearpage

\section{Pseudo modo deslizante}

Considérese la superficie de deslizamiento
\[
  \boxed{\sigma=X_2+mX_1},
  \qquad m>0.
\]

En lugar de utilizar una ley discontinua, se propone
\[
  \boxed{u=-mX_2-K\operatorname{sat}(\sigma)}.
\]

La función de saturación se define como
\[
  \operatorname{sat}(X)=
  \begin{cases}
    -L, & X<-L,\\
    X, & |X|\leq L,\\
    L, & X>L,
  \end{cases}
\]
donde \(L>0\). También es posible utilizar una función sigmoide, por ejemplo
\[
  \tanh(X).
\]

La pendiente de estas funciones puede modificarse mediante un parámetro
\(\eta>0\):
\[
  \operatorname{sat}\left(\frac{X}{\eta}\right),
  \qquad
  \tanh\left(\frac{X}{\eta}\right).
\]

\subsection{Aplicación al modelo PVTOL}

Después de la linealización exacta por realimentación, se tienen las
dinámicas desacopladas
\[
  \ddot{x}=v_x,\qquad
  \ddot{z}=v_z,\qquad
  \ddot{\theta}=\tau.
\]

Se conservan los errores utilizados en el controlador anterior:
\[
  e_x=r_x-x,\qquad
  e_z=r_z-z,\qquad
  e_\theta=r_\theta-\theta,
\]
y sus derivadas:
\[
  \dot e_x=\dot r_x-\dot x,\qquad
  \dot e_z=\dot r_z-\dot z,\qquad
  \dot e_\theta=\dot r_\theta-\dot\theta.
\]

Las superficies de pseudo modo deslizante son
\begin{align*}
  \sigma_x&=\dot e_x+m_xe_x,\\
  \sigma_z&=\dot e_z+m_ze_z,\\
  \sigma_\theta&=\dot e_\theta+m_\theta e_\theta.
\end{align*}

Para implementar la saturación general con amplitud \(L\), se utiliza
\[
  \operatorname{sat}\left(\frac{X}{\eta}\right)
  =
  \begin{cases}
    -L, & X<-L\eta,\\
    \dfrac{X}{\eta}, & |X|\leq L\eta,\\
    L, & X>L\eta.
  \end{cases}
\]

La ley aplicada al PVTOL conserva el signo de realimentación de las entradas
virtuales empleadas en el modelo:
\begin{align*}
  v_x&=-m_x\dot x
  +K_x\operatorname{sat}\left(\frac{\sigma_x}{\eta_x}\right),\\
  v_z&=-m_z\dot z
  +K_z\operatorname{sat}\left(\frac{\sigma_z}{\eta_z}\right),\\
  \tau&=-m_\theta\dot\theta
  +K_\theta\operatorname{sat}\left(\frac{\sigma_\theta}{\eta_\theta}\right).
\end{align*}

En cada canal se pueden utilizar amplitudes independientes:
\[
  L_x>0,\qquad L_z>0,\qquad L_\theta>0.
\]
Así, las saturaciones implementadas son
\[
  \operatorname{sat}_x\left(\frac{\sigma_x}{\eta_x}\right)
  =
  \max\left(-L_x,\min\left(L_x,\frac{\sigma_x}{\eta_x}\right)\right),
\]
\[
  \operatorname{sat}_z\left(\frac{\sigma_z}{\eta_z}\right)
  =
  \max\left(-L_z,\min\left(L_z,\frac{\sigma_z}{\eta_z}\right)\right),
\]
\[
  \operatorname{sat}_\theta\left(\frac{\sigma_\theta}{\eta_\theta}\right)
  =
  \max\left(-L_\theta,
  \min\left(L_\theta,\frac{\sigma_\theta}{\eta_\theta}\right)\right).
\]

La referencia angular continúa siendo
\[
  r_\theta:=\theta_d
  =\operatorname{atan2}(-v_x,v_z+g),
\]
y las entradas físicas se obtienen mediante
\[
  \boxed{U_d=\sqrt{v_x^2+(v_z+g)^2}}.
\]
